\chapter{Stochastic Diffusion}\label{ch:diffusion}

The resolved velocity field carries particles by advection, but it omits the mixing done by motions below
the model's resolution. The engine can represent this unresolved mixing as a random walk superimposed on
the advective step, turning the trajectory equation into a stochastic differential equation. This chapter
sets out the random-walk model and its treatment of the integration direction, how a horizontal kick moves
a particle on the sphere, the three horizontal schemes, the way the diffusivity and its gradient are
sampled from the mesh, the vertical random walk, the counter-based random-number generator that makes the
result reproducible and restart-safe, and how the kicks compose with the integrators. Diffusion is off by
default; when disabled it is byte-for-byte identical to a purely advective run.

\section{The random-walk model}\label{sec:diff-model}

A diffusivity $K_h$ spreads an ensemble of particles so that its variance grows linearly in time,
$\sigma^2 = 2K_h t$. The corresponding increment added to a particle over a step combines the
deterministic advective displacement with a random kick (\cref{fig:diffusion}):
\begin{equation}\label{eq:diff-step}
  \Delta x = u\,\dt \;+\; \xi\,\sqrt{2K_h\,\lvert\dt\rvert},
\end{equation}
where $\xi$ is a standard normal random variate drawn independently per component. With a spatially varying
diffusivity a drift term $\partial_x K_h\,\lvert\dt\rvert$ joins them (\cref{sec:diff-schemes}). The signs
follow three rules. The advective term carries the signed step, so it reverses under backward integration.
The random term uses $\sqrt{\lvert\dt\rvert}$ and spreads the ensemble in \emph{either} direction. The
drift uses $\lvert\dt\rvert$ too: it belongs to the diffusion, not to the advection, so a backward run
integrates the adjoint process, $-u+\nabla K_h$ plus the same noise. (A drift that reversed with the step
would leave a backward ensemble in a closed column with a stationary density proportional to $K^{-2}$
instead of a uniform one.) Diffusion is intrinsically irreversible; a backward run does not un-mix what a
forward run mixed, and the backward cloud is not a validated probability distribution of source positions.
The diffusivity may be anisotropic, with independent zonal and meridional components $K_x,K_y$.

\begin{figure}[t]
\centering
\begin{tikzpicture}[scale=1.15]
  \node[meshnode,label={[font=\footnotesize]below:$x_n$}] (A) at (0,0) {};
  \draw[-{Stealth[length=6pt]},lagblue,line width=1.2pt] (A)--(2.6,0.7) node[midway,above,font=\footnotesize,lagblue!70!black]{$u\,\dt$};
  \coordinate (B) at (2.6,0.7);
  \node[circle,fill=lagblue,inner sep=1.2pt] at (B) {};
  \foreach \a in {20,70,130,200,255,310,350,160,95}
     \draw[densely dotted,kelp!75] (B)--($(B)+(\a:0.62)$);
  \draw[dashed,kelp,line width=0.8pt] (B) circle (0.62);
  \node[font=\footnotesize,kelp!70!black] at (3.95,0.7) {$\xi\sqrt{2K_h|\dt|}$};
  \foreach \a in {30,110,240,330} \node[drifter,scale=0.5] at ($(B)+(\a:0.62)$) {};
  \node[font=\footnotesize,gray] at (1.0,-0.55) {deterministic};
  \node[font=\footnotesize,gray] at (3.0,-0.45) {stochastic kick};
\end{tikzpicture}
\caption{One diffusive step. The deterministic advective displacement $u\,\dt$ (blue) carries the particle
from $x_n$; a random kick of typical size $\sqrt{2K_h|\dt|}$ (dashed circle) is then added, scattering the
realized endpoint (stars). The advective term reverses under backward integration; the kick, scaling with
$\sqrt{|\dt|}$, spreads the ensemble either way.}
\label{fig:diffusion}
\end{figure}

\section{Kicks on the sphere}\label{sec:diff-sphere}

A kick is a displacement $(d_x,d_y)$ in metres, east and north. Added to longitude and latitude through the
metric of \cref{eq:int-displacement}, it would be a random walk on the longitude--latitude chart, which lacks
the drift $-K\tan\latx/R_M$ that a random walk on the sphere has: its particles would spread uniformly in
latitude rather than in area and would accumulate toward the poles. The engine therefore moves every
horizontal kick along the great circle. With $d=\sqrt{d_x^2+d_y^2}$, $\delta=d/R_M$ and the metric radius $R_M=\Mlat\cdot180/\pi$,
the position $\mathbf{r}=(\cos\latx\cos\lonx,\cos\latx\sin\lonx,\sin\latx)$ and the tangent displacement
\begin{equation}\label{eq:diff-tangent}
  \mathbf{t}=d_x\,(-\sin\lonx,\ \cos\lonx,\ 0)+d_y\,(-\sin\latx\cos\lonx,\ -\sin\latx\sin\lonx,\ \cos\latx),
\end{equation}
the kicked position is
\begin{equation}\label{eq:diff-expmap}
  \mathbf{r}'=\mathbf{r}\cos\delta+\frac{\mathbf{t}}{d}\sin\delta ,
\end{equation}
converted back to longitude and latitude. The east and north directions are taken at the stored longitude,
so the map is regular through a pole. For an anisotropic diffusivity the sphere adds one more term, a
northward drift $(K_x-K_y)\tan\latx\,\lvert\dt\rvert/R_M$ (with the latitude capped at $\pm89.9^\circ$); with
$K_x=K_y$ it vanishes and the uniform scheme is free of any drift.

\codeanchor{\fn{sphere\_displace}, \fn{sphere\_displace\_arrays} and \fn{anisotropic\_metric\_drift\_m} in
\srcfile{numerics/diffusion.py}.}

\section{The three horizontal schemes}\label{sec:diff-schemes}

The key \conf{diffusion.scheme} selects among three realizations of increasing fidelity to a spatially
varying diffusivity. The \conf{uniform} scheme uses a constant diffusivity taken from the configuration
(\conf{kh\_zonal}, \conf{kh\_meridional}) and adds the kick of \cref{eq:diff-step} with no gradient term.
The \conf{em} scheme (Euler--Maruyama) reads a spatially varying diffusivity from the dataset and adds the
Itô drift, per component,
\begin{equation}\label{eq:diff-em}
  \Delta x = u\,\dt+\partial_x K_x\,\lvert\dt\rvert + \sqrt{2K_x\lvert\dt\rvert}\;\xi_1 ,
\end{equation}
the drift being what prevents particles from accumulating spuriously in regions of low diffusivity and so
keeps the scheme consistent with the diffusion equation. The \conf{m1} scheme (Milstein) adds the next
term, $\tfrac12\partial_x K_x\,\lvert\dt\rvert\,(\xi_1^2-1)$; it is rigorous only for a separable diagonal
diffusivity and warns once whenever the field varies in space, for which \conf{em} is the safer
choice. With Euler (the combined composition) the diffusive part --- drift, Milstein term and noise ---
moves the particle along the great circle from the start of the step (\cref{sec:diff-sphere}) and the
advective part is added on the chart; with RK2 and RK4 (the split composition) the kick is applied at the
position the deterministic step reached. Where the diffusivity is constant, \conf{em} and \conf{m1} draw
the same normals as \conf{uniform}; only where the kick is applied differs.

\section{Sampling the diffusivity}\label{sec:diff-sampling}

The spatially varying schemes require both the diffusivity and its spatial gradient at the particle. The
diffusivity is read from node or element variables, their layout detected from the dimensions, and it must
not be negative. The value is obtained by the least-squares plane fit of \cref{sec:hvel-leastsq}, gated on
containment. The gradient is obtained by central differencing of that value over an offset
\conf{dres\_m} --- by default the median edge of up to \num{5000} sampled elements, clipped to
\SIrange{100}{50000}{\metre} --- with a fallback chain to a one-sided difference and then to the analytic
gradient of the fitted plane should the central difference be unavailable near a boundary.

\section{Vertical diffusion}\label{sec:diff-vertical}

In three dimensions on a z-level or sigma mesh, a vertical diffusivity $A_z$ adds a vertical random walk
(\conf{az\_enabled}, with a constant \conf{az\_const} or, on FESOM2 z-level output, a field \conf{var\_az}
on the interfaces). Its increment in positive-down depth is
\begin{equation}\label{eq:diff-az}
  \Delta\depx=\frac{\partial A_z}{\partial\depx}\,\lvert\dt\rvert+\sqrt{2A_z\lvert\dt\rvert}\;\xi ,
\end{equation}
with the drift mandatory: without it a walk in a non-uniform $A_z$ would collect particles where $A_z$ is
small \citep{visser1997random}. The field is taken from the nearest node's column, linear between
interfaces, with the depth clamped to the wet part of the column; the gradient is the slope of the same
linear segment --- on an interface, the segment above --- so that the value that sets the noise and the
gradient that sets the drift belong to one interpolant. A kick that leaves the column is reflected, not
clamped: with the column $[a,b]$, $L=b-a$ and $y=(\depx-a)\bmod 2L$, the new depth is $a+y$ for $y\le L$
and $a+2L-y$ otherwise, which handles any number of bounces. A clamp would pile particles up at the
boundaries; the reflection keeps a well-mixed column well mixed, which the tests check in a diffusivity
that varies tenfold over the column, in both directions of integration. The column is the surface to the
nearest node's floor on a z-level mesh and the particle's end-of-step column less $z_{0b}$ on a sigma mesh,
the same bounds as the depth clamp of \cref{sec:vert-clamp}. The vertical kick composes only with the
\conf{uniform} horizontal scheme, runs after the deterministic step and the horizontal kick, and with no
source configured the facility is dormant.

\codeanchor{\fn{diffusion\_az\_step\_meters}, \fn{az\_column\_value\_gradient}, \fn{reflect\_depth} and the
loops \fn{\_loop\_az\_kick\_zlevel} and \fn{\_loop\_az\_kick\_sigma} (\srcfile{numerics/diffusion.py});
applied by \fn{\_apply\_split\_az\_kick} (\srcfile{numerics/integrators.py}).}

\section{Reproducibility: counter-based random numbers}\label{sec:diff-rng}

The random kicks are generated not by a sequential pseudo-random stream but by a \emph{counter-based}
generator \citep{salmon2011philox}, in which each draw is a stateless hash of an integer key. The key is the
tuple of the configuration seed, the realization index, the \emph{original} particle index, the global step
index, and the component --- zero for the horizontal pair, one for the vertical kick --- from which a 64-bit
hash produces a uniform pair that is converted to a normal pair. Because each draw depends only on its key
and not on the order in which draws are issued, the result is bit-for-bit identical regardless of the
number of threads or the order of execution, and it is identical across a restart --- the same step index
regenerates the same kick. The particle index in the key must be the original, fixed index, not a
compacted index that shifts as particles die, or the stream would change when the live set changes. A run
over one shard of a larger release (\cref{sec:perf-shard}) adds the global index of its first particle,
\conf{data.particle\_offset} or the header that the splitting tool writes into each shard file, so a
particle draws the same numbers in the shard as in the unsharded run.

\section{Composition with the integrator}\label{sec:diff-composition}

How the stochastic kick combines with the deterministic step depends on the scheme and the integrator
(\conf{diffusion.composition}, resolved when the run starts). The \conf{uniform} scheme composes by
\emph{splitting}: the deterministic step (of any integrator, including RK45) is taken first and the kick
added afterward. The \conf{em} and \conf{m1} schemes with the Euler integrator \emph{combine}: the
stochastic differential equation replaces the Euler step as a whole; this mode needs element-layout
least-squares velocities. The same schemes with RK2 or RK4 use an opt-in split
(\conf{composition: split\_rk}), taking the Runge--Kutta step and then the kick, with the diffusivity
sampled at the start or, by default, at the end of the step (\conf{kh\_eval\_position}) and the kick applied
at the post-step position; the engine warns once that this is a pragmatic extension rather than a rigorous
scheme. The spatially varying schemes do not combine with RK45.

The kicked position is checked like any other: a kick that leaves the mesh returns the particle to its
position at the \emph{start} of the step with status \code{OUTSIDE}. In three dimensions only the
\conf{uniform}/split combination is available for the horizontal kick; the depth is clamped again at the
kicked position, and on a sigma mesh the particle keeps its fraction of the column. The combined and
split-with-Runge--Kutta modes raise an explicit error in three dimensions rather than apply an unvalidated
scheme. In time, the diffusivities are interpolated linearly with the velocity payload and are not included
in the monotone-cubic reconstruction.

\codeanchor{The schemes, the drift, and the Milstein term are \fn{diffusion\_uniform\_step\_meters},
\fn{diffusion\_em\_step\_meters} and \fn{diffusion\_m1\_step\_meters}; the counter-based generator --- a
keyed SplitMix64 hash in the counter-based spirit of \citet{salmon2011philox} --- and the Box--Muller pair
are in the same module, \srcfile{numerics/diffusion.py}. The composition is resolved by
\fn{\_resolve\_composition} (\srcfile{tracking.py}) and applied by \fn{\_apply\_split\_diffusion\_kick},
\fn{\_apply\_combined\_em\_step} and \fn{\_apply\_split\_rk\_em\_kick} (\srcfile{numerics/integrators.py});
the default offset is \fn{\_compute\_default\_dres\_m} (\srcfile{tracking.py}).}
